(1) 阱内波函数满足 \(\psi(0)=\psi(L)=0\)，归一化基态族为 \(\psi_n(x)=\sqrt{2/L}\sin(n\pi x/L)\)，\(n=1,2,\ldots\)。能量 \(E_n^{(0)}=n^2\pi^2\hbar^2/(2mL^2)\)，各能级在一维中非简并；中心 \(\delta\) 势是有限积分强度的短程微扰。
(2) 一阶非简并微扰为 \(\Delta E_n^{(1)}=\langle n|V_1|n\rangle=\lambda|\psi_n(L/2)|^2=(2\lambda/L)\sin^2(n\pi/2)\)。当 \(n\) 为奇数，中心位移波函数不为零，线性修正为 \(2\lambda/L\)；当 \(n\) 为偶数，中心恰为节点，一阶修正为零。
(3) 对偶数 \(n=2j\)，有 \(\psi_{2j}(L/2)=\sin(j\pi)=0\)。中心\(\delta \)势对该态的作用 \(V_1\psi_{2j}=\lambda\psi_{2j}(L/2)\delta(x-L/2)=0\)，且波函数和导数在中心连续，故它严格满足加势后的薛定谔方程。本征能 \(E_{2j}=E_{2j}^{(0)}\) 对任意有限 \(\lambda\) 都不变，而不只是某一阶微扰的巧合。
(4) 当 \(\lambda\to+\infty\)，中心位置几乎成为不可穿越节点，原阱近似分裂为两个宽度 \(L/2\) 的独立无限方阱。每个半阱最低能量为 \(\pi^2\hbar^2/[2m(L/2)^2]=4E_1^{(0)}=E_2^{(0)}\)；原偶数节点态从始至终位于这些能量，原奇数能级随屏障变高上移并逼近相邻偶数能级，形成左右半阱的近简并组合。
